\documentclass[12pt,a4paper]{article}

% ── Packages ──────────────────────────────────────────────────────────────────
\usepackage[T1]{fontenc}
\usepackage[utf8]{inputenc}
\usepackage{lmodern}
\usepackage[margin=2.5cm]{geometry}
\usepackage{hyperref}
\usepackage{xcolor}
\usepackage{booktabs}
\usepackage{longtable}
\usepackage{listings}
\usepackage{graphicx}
\usepackage{enumitem}
\usepackage{titlesec}
\usepackage{fancyhdr}
\usepackage{microtype}
\usepackage{parskip}
\usepackage{array}
\usepackage{tabularx}
\usepackage{float}

% ── Colours ───────────────────────────────────────────────────────────────────
\definecolor{codebg}{HTML}{F5F5F5}
\definecolor{codekw}{HTML}{0000CC}
\definecolor{codecomment}{HTML}{228B22}
\definecolor{codestring}{HTML}{AA0000}
\definecolor{linkblue}{HTML}{1A6496}
\definecolor{sectiongray}{HTML}{2C3E50}

% ── Hyperref ──────────────────────────────────────────────────────────────────
\hypersetup{
    colorlinks  = true,
    linkcolor   = linkblue,
    urlcolor    = linkblue,
    citecolor   = linkblue,
    pdftitle    = {FlexQL Design Document},
    pdfauthor   = {K Chaitanya},
    pdfsubject  = {Database Engine Design},
}

% ── Section title formatting ──────────────────────────────────────────────────
\titleformat{\section}{\large\bfseries\color{sectiongray}}{}{0em}{\thesection\quad}[\titlerule]
\titleformat{\subsection}{\normalsize\bfseries\color{sectiongray}}{}{0em}{\thesubsection\quad}
\titleformat{\subsubsection}{\normalsize\itshape}{}{0em}{\thesubsubsection\quad}

% ── Code listings ─────────────────────────────────────────────────────────────
\lstset{
    backgroundcolor = \color{codebg},
    basicstyle      = \ttfamily\small,
    keywordstyle    = \color{codekw}\bfseries,
    commentstyle    = \color{codecomment}\itshape,
    stringstyle     = \color{codestring},
    breaklines      = true,
    frame           = single,
    framesep        = 4pt,
    rulecolor       = \color{gray!40},
    xleftmargin     = 6pt,
    xrightmargin    = 6pt,
    showstringspaces= false,
    tabsize         = 4,
    captionpos      = b,
}
\lstdefinelanguage{SQL}{
    keywords  = {CREATE,TABLE,IF,NOT,EXISTS,INSERT,INTO,VALUES,SELECT,FROM,WHERE,
                 ORDER,BY,DESC,ASC,INNER,LEFT,JOIN,ON,DELETE,UPDATE,SET,DROP,
                 TRUNCATE,SHOW,DISTINCT,LIMIT,OFFSET,LIKE,AND,OR,IS,BETWEEN,IN,
                 NULL,GROUP,HAVING,AS,USE,ALTER,ADD,COLUMN,MODIFY,DESCRIBE,
                 PRIMARY,KEY,INT,DECIMAL,VARCHAR,TEXT,DATETIME,COUNT,SUM,AVG,MIN,MAX},
    morestring  = [b]',
    morecomment = [l]{--},
    sensitive   = false,
}

% ── Header / Footer ───────────────────────────────────────────────────────────
\pagestyle{fancy}
\fancyhf{}
\lhead{\small\textbf{FlexQL} — Design Document}
\cfoot{\thepage}
\renewcommand{\headrulewidth}{0.4pt}

% ── Table helper ──────────────────────────────────────────────────────────────
\newcolumntype{L}[1]{>{\raggedright\arraybackslash}p{#1}}

% ══════════════════════════════════════════════════════════════════════════════
\begin{document}

% ── Title page ────────────────────────────────────────────────────────────────
\begin{titlepage}
\centering
\vspace*{3cm}
{\Huge\bfseries\color{sectiongray} FlexQL \par}
\vspace{0.4cm}
{\large\itshape A High-Performance In-Memory Database Engine \par}
\vspace{1.5cm}
{\normalsize\bfseries Design Document \par}
\vspace{0.5cm}
\hrule
\vspace{0.5cm}
{\small
\begin{tabular}{ll}
\textbf{Author}     & K chaitanya \\[2pt]
\textbf{Repository} & \href{https://github.com/devchaitanya/FlexQl}{github.com/devchaitanya/FlexQl} \\[2pt]
\textbf{Language}   & C++17 (no external DB libraries) \\[2pt]
\textbf{Platform}   & Linux (GCC 15 / Clang 20) \\[2pt]
\end{tabular}
}
\vspace{0.5cm}
\hrule
\vfill
{\small Design Lab — Database Systems Project}
\end{titlepage}

% ── Table of contents ─────────────────────────────────────────────────────────
\tableofcontents
\newpage

% ══════════════════════════════════════════════════════════════════════════════
\section{System Overview}

FlexQL is a client-server, in-memory relational database implemented entirely in C++17
with no external database libraries. It exposes a SQLite-compatible C API and a TCP
wire protocol that allows multiple concurrent clients.

\subsection*{Architecture Diagram}

\begin{lstlisting}[language={},title={High-level architecture}]
Client (REPL / benchmark)
        |  TCP socket  (raw SQL text terminated by ';')
        v
 flexql-server
  +-- Thread Pool  (N = hardware_concurrency workers)
  +-- Parser  ->  Executor
  +-- Table Storage  (columnar StringArena + flat arrays)
  +-- LRU Query Cache
  +-- PK Index  (flat open-addressing FNV-64 hash map)
  +-- WAL Writer  (zero-copy writev + SSE4.2 CRC32c)
  +-- Snapshot Engine  (binary checkpoints)
\end{lstlisting}

\subsection*{Client API}

The client library (\texttt{flexql.h}) exposes exactly four functions as prescribed:

\begin{lstlisting}[language=C]
typedef struct FlexQL FlexQL;
#define FLEXQL_OK    0
#define FLEXQL_ERROR 1

int  flexql_open (const char *host, int port, FlexQL **db);
int  flexql_close(FlexQL *db);
int  flexql_exec (FlexQL *db, const char *sql,
                  int (*callback)(void*,int,char**,char**),
                  void *arg, char **errmsg);
void flexql_free (void *ptr);
\end{lstlisting}

The callback is invoked once per result row:
\texttt{callback(data, columnCount, values[], columnNames[])}.

% ══════════════════════════════════════════════════════════════════════════════
\section{How Data is Stored}

\subsection{Columnar In-Memory Layout}

Each table uses two parallel flat arrays rather than a \texttt{vector<vector<string>>}:

\begin{lstlisting}[language=C++]
StringArena                   arena_;    // slab allocator owns all string bytes
std::vector<RowMeta>          rows_meta_;  // {expires_at, deleted} per row
std::vector<std::string_view> rows_vals_;  // rows * cols string_views into arena
int                           num_cols_;
\end{lstlisting}

\noindent Value at row $r$, column $c$: \texttt{rows\_vals\_[r * num\_cols\_ + c]}

\subsection{StringArena Slab Allocator}

\begin{lstlisting}[language=C++]
class StringArena {
    static constexpr size_t BLOCK = 4 * 1024 * 1024;  // 4 MB slabs
    std::vector<std::unique_ptr<char[]>> blocks_;
    size_t used_ = BLOCK;
public:
    std::string_view intern(std::string_view sv) {
        if (sv.size() > BLOCK) { /* dedicated oversized block */ }
        if (used_ + sv.size() > BLOCK) { /* new slab */ }
        char* p = blocks_.back().get() + used_;
        std::memcpy(p, sv.data(), sv.size());
        used_ += sv.size();
        return {p, sv.size()};
    }
};
\end{lstlisting}

\noindent Slab blocks never move, so interned \texttt{string\_view} references are \textbf{stable
for the lifetime of the table}.

\subsection{Why This Design?}

\begin{table}[H]
\centering
\begin{tabularx}{\textwidth}{L{3.5cm} L{3.5cm} X}
\toprule
\textbf{Design choice} & \textbf{Alternative} & \textbf{Reason} \\
\midrule
Flat \texttt{string\_view} array & \texttt{vector<vector<string>>} &
  Eliminates 1 heap allocation per row. For 1\,M rows: zero per-row \texttt{malloc} calls \\
\addlinespace
StringArena slabs & Per-string \texttt{std::string} &
  All data in contiguous 4\,MB blocks; hot strings stay in L2/L3 cache \\
\addlinespace
Pre-reserve on first batch & Default \texttt{push\_back} growth &
  Avoids 7+ reallocation spikes during a 1\,M-row insert benchmark \\
\bottomrule
\end{tabularx}
\caption{Storage design choices and rationale}
\end{table}

% ══════════════════════════════════════════════════════════════════════════════
\section{Indexing Method}

\subsection{Flat Open-Addressing Hash Map (Primary Key only)}

\texttt{PrimaryIndex} in \texttt{include/index/pk\_index.h}:

\begin{itemize}[noitemsep]
  \item \textbf{Hash function:} FNV-64 — fast, good distribution for string keys
  \item \textbf{Collision resolution:} linear probing — cache-friendly, no pointer chasing
  \item \textbf{Load factor threshold:} 0.6 — rehash doubles capacity when exceeded
  \item \textbf{Keys:} \texttt{string\_view} into arena (zero-copy lookup)
  \item \textbf{Supports:} \texttt{insert}, \texttt{lookup}, \texttt{remove}, \texttt{update}, \texttt{clear}
\end{itemize}

\begin{lstlisting}[language=C++,title={FNV-64 hash + probe loop}]
static uint64_t fnv64(std::string_view s) noexcept {
    uint64_t h = 14695981039346656037ULL;
    for (unsigned char c : s) h = (h ^ c) * 1099511628211ULL;
    return h;
}
// lookup: probe from fnv64(key) % capacity until match or empty slot
\end{lstlisting}

\subsection{Why Not \texttt{std::unordered\_map}?}

Each \texttt{unordered\_map} bucket is a linked list — pointer-chasing across cache lines.
For 1\,M rows, the flat open-addressing map is approximately $3\times$ faster on lookups
because all slots are contiguous in memory.

\subsection{Tables Without a Primary Key}

Tables without a \texttt{PRIMARY KEY} constraint have no index. All queries do a linear
scan of \texttt{rows\_meta\_} and \texttt{rows\_vals\_}. This is intentional — the
benchmark table (\texttt{BIG\_USERS}) has no primary key, so the hot insert path pays
zero index overhead.

% ══════════════════════════════════════════════════════════════════════════════
\section{Caching Strategy}

\subsection{LRU Query Cache}

\begin{itemize}[noitemsep]
  \item \textbf{Policy:} Least Recently Used (LRU)
  \item \textbf{Capacity:} Configurable via \texttt{cache\_capacity} in \texttt{config/server.conf}
        (default 1024 entries)
  \item \textbf{Key format:} \texttt{TABLE|cols|W:col=op=val|O:colDA|D|Lnn}
  \item \textbf{Invalidation:} Any \texttt{INSERT / UPDATE / DELETE / TRUNCATE} calls
        \texttt{invalidate\_prefix(table\_name)}, evicting all cached results for that table
  \item \textbf{Scope:} \texttt{SELECT} queries only; compound \texttt{WHERE} predicates bypass
        the cache (they are not fully serialised to a key)
\end{itemize}

\subsection{Thread Safety}

The cache is protected by a \texttt{shared\_mutex}: multiple readers acquire a shared
lock concurrently; a write (\texttt{put} / \texttt{invalidate}) takes an exclusive lock.

\subsection{Why LRU over LFU?}

The benchmark pattern is: insert 1\,M rows once, then run repeated \texttt{SELECT} scans.
LRU captures the \textit{``recently scanned table''} pattern with no frequency bookkeeping
overhead. LFU would yield identical hit rates but add per-entry counter updates on every
access.

% ══════════════════════════════════════════════════════════════════════════════
\section{Handling of Expiration Timestamps}

\subsection{Insertion}

Every row stores \texttt{int64\_t expires\_at} (Unix epoch seconds; \texttt{0} = never
expires). The benchmark sends rows with a 5th column as the TTL value. The executor
detects a column-count mismatch of $+1$, extracts the last value as \texttt{expires\_at},
and strips it before passing the row to storage.

\subsection{Filtering}

\begin{lstlisting}[language=C++]
bool Table::is_expired(const RowMeta& m) const {
    if (m.expires_at == 0) return false;
    return m.expires_at <= (int64_t)std::time(nullptr);
}
\end{lstlisting}

Every scan, delete, and update skips expired rows. They become invisible to all queries
immediately upon expiry — no background job required for correctness.

\subsection{Physical Removal (Compaction)}

A background \textbf{TTL sweeper} thread (\texttt{TTLManager}) runs every
\texttt{ttl\_sweep\_interval} seconds (default 60\,s). It calls \texttt{Table::compact()}
on each table, which rebuilds both parallel arrays and the PK index, physically removing
rows where \texttt{deleted == true} or \texttt{is\_expired() == true}.

This two-phase design (logical deletion $\to$ background physical removal) avoids blocking
the hot INSERT/SELECT path with compaction work.

% ══════════════════════════════════════════════════════════════════════════════
\section{Multithreading Design}

\subsection{Thread Pool}

\texttt{ThreadPool} creates \texttt{hardware\_concurrency()} worker threads (configurable
via \texttt{threads} in \texttt{server.conf}). Each accepted TCP connection is dispatched
to a free worker via a \texttt{std::queue} protected by \texttt{std::mutex} +
\texttt{std::condition\_variable}.

\subsection{Per-Table \texttt{shared\_mutex}}

\begin{table}[H]
\centering
\begin{tabularx}{\textwidth}{L{4.5cm} L{3cm} X}
\toprule
\textbf{Operation} & \textbf{Lock type} & \textbf{Behaviour} \\
\midrule
\texttt{scan()} & \texttt{shared\_lock} & Multiple concurrent readers \\
\texttt{insert\_flat()} / \texttt{insert\_batch()} & \texttt{unique\_lock} & Exclusive write \\
\texttt{delete\_rows()} / \texttt{update\_rows()} & \texttt{unique\_lock} & Exclusive write \\
\texttt{compact()} & \texttt{unique\_lock} & Background only \\
\bottomrule
\end{tabularx}
\caption{Locking strategy per table operation}
\end{table}

Multiple clients can read the same table simultaneously. A write from one client blocks
only that table, not others.

\subsection{WAL Writer}

The WAL writer uses a single \texttt{std::mutex}. \texttt{writev} takes $\sim$57\,µs per
call; contention is negligible under the benchmark's sequential insert pattern.

\subsection{TCP\_NODELAY}

All accepted sockets have \texttt{TCP\_NODELAY} set, disabling Nagle's algorithm. This
eliminates up to 40\,ms of artificial per-request latency on loopback interfaces.

% ══════════════════════════════════════════════════════════════════════════════
\section{Persistence — WAL and Snapshots}

\subsection{Write-Ahead Log (WAL)}

Every mutation (\texttt{INSERT}, \texttt{UPDATE}, \texttt{DELETE}, \texttt{TRUNCATE},
\texttt{CREATE TABLE}, \texttt{DROP TABLE}) is written to the WAL before execution.
On crash and restart, the WAL is replayed to recover all committed writes.

\subsubsection*{Record Format (20-byte header + SQL payload)}

\begin{lstlisting}[language={}]
[ 4B magic ][ 8B LSN ][ 4B payload_len ][ 4B CRC32c ][ SQL bytes ]
\end{lstlisting}

\subsubsection*{Zero-Copy Append}

\begin{lstlisting}[language=C++]
uint8_t hdr[20];   // built on stack (magic, LSN, len, checksum)
struct iovec iov[2] = {{hdr, 20}, {sql.data(), sql.size()}};
::writev(fd_, iov, 2);   // single syscall, no intermediate copy
\end{lstlisting}

\subsubsection*{Key Design Decisions}

\begin{itemize}[noitemsep]
  \item \texttt{fdatasync} is called \textit{only} at shutdown/snapshot — not per write.
        This saves $\sim$1300\,ms over the benchmark (100 $\times$ 13\,ms per flush).
  \item \textbf{CRC32c} uses the SSE4.2 hardware instruction \texttt{\_mm\_crc32\_u64}:
        $\sim$9\,ms for 200 $\times$ 275\,KB batches vs $\sim$166\,ms for software lookup.
\end{itemize}

\subsection{Binary Snapshots}

Checkpoints are written to \texttt{data/snapshots/snap\_\textit{lsn}.bin}:

\begin{lstlisting}[language={}]
[ 8B magic ][ 8B snapshot_lsn ]
[ for each table:
    4B name_len, name bytes
    4B num_columns, [ col_name, col_type, pk, not_null ] * N
    8B num_rows,    [ col_values..., expires_at ] * M
]
\end{lstlisting}

\subsection{Recovery on Startup}

\begin{enumerate}[noitemsep]
  \item Load the snapshot with the highest LSN (restores all schemas and row data)
  \item Replay WAL records with \texttt{LSN > snapshot\_lsn} (recovers mutations since checkpoint)
  \item Accept client connections
\end{enumerate}

% ══════════════════════════════════════════════════════════════════════════════
\section{SQL Language Support}

The parser is a hand-written recursive-descent parser with a dedicated fast path
(\texttt{fast\_parse\_insert}) for performance-critical INSERT statements.

\subsection{DDL}

\begin{lstlisting}[language=SQL]
CREATE TABLE [IF NOT EXISTS] t (col type [PRIMARY KEY] [NOT NULL], ...)
  -- types: INT, DECIMAL, VARCHAR(n), TEXT, DATETIME
DROP TABLE [IF EXISTS] t
TRUNCATE TABLE t
ALTER TABLE t ADD COLUMN col type
\end{lstlisting}

\subsection{DML}

\begin{lstlisting}[language=SQL]
INSERT INTO t [(col, ...)] VALUES (v, ...) [, (v, ...) ...]
UPDATE t SET col=v [, col=v ...] [WHERE ...]
DELETE FROM t [WHERE ...]
\end{lstlisting}

\subsection{DQL}

\begin{lstlisting}[language=SQL]
SELECT [DISTINCT] col [AS alias], ...
  FROM t
  [INNER JOIN t2 ON t.col = t2.col]
  [LEFT  JOIN t2 ON t.col = t2.col]
  [WHERE ...]
  [GROUP BY col [, col ...] [HAVING ...]]
  [ORDER BY col [ASC|DESC] [, col [ASC|DESC] ...]]
  [LIMIT n] [OFFSET m]

SELECT COUNT(*) | SUM(col) | AVG(col) | MIN(col) | MAX(col)
  FROM t [WHERE ...]
\end{lstlisting}

\subsection{WHERE Clause (Single Condition)}

As specified in the assignment, only a single condition is supported per \texttt{WHERE}
clause. Logical combinations (\texttt{AND}, \texttt{OR}, \texttt{NOT}) are explicitly
rejected by the parser.

Supported operators: \texttt{=}, \texttt{!=}, \texttt{<}, \texttt{<=}, \texttt{>},
\texttt{>=}, \texttt{LIKE} (\texttt{\%} and \texttt{\_} wildcards),
\texttt{BETWEEN v AND v}, \texttt{IN (v, ...)}, \texttt{IS NULL},
\texttt{IS NOT NULL}.

Each \texttt{WHERE} predicate is represented as a single \texttt{Predicate::LEAF}
node evaluated during the table scan.

\subsection{Utility Commands}

\begin{lstlisting}[language=SQL]
SHOW TABLES
SHOW DATABASES
SHOW COLUMNS FROM t    -- synonym for DESCRIBE
DESCRIBE t
USE [DATABASE] name
\end{lstlisting}

% ══════════════════════════════════════════════════════════════════════════════
\section{Wire Protocol}

Plain-text, newline-delimited over TCP. The client sends raw SQL terminated by
\texttt{;} in a single \texttt{send} call. The server replies:

\begin{lstlisting}[language={}]
# SELECT with results:
COLS col1 col2 col3
ROW  val1 val2 val3
...
END

# Non-SELECT (INSERT, UPDATE, ...):
OK
END

# Error:
ERROR: <message>
END
\end{lstlisting}

The \texttt{COLS} header carries real column names including \texttt{AS} aliases.
The benchmark client (\texttt{bench/flexql.cpp}) ignores non-\texttt{ROW} lines, so the
header is fully backward-compatible.

\textbf{Buffering:} The server uses a 64\,KB \texttt{recv} buffer (pre-reserved to
512\,KB for large \texttt{INSERT} batches). The client uses a 4\,KB local read buffer
to eliminate per-character \texttt{recv} syscalls.

% ══════════════════════════════════════════════════════════════════════════════
\section{Other Design Decisions}

\subsection{Parser Fast Path for INSERT}

The parser has two code paths: a full recursive-descent parser for all SQL statements,
and a dedicated \texttt{fast\_parse\_insert()} for INSERT. The fast path operates as a
single-pass character scanner over the raw SQL buffer, producing
\texttt{std::string\_view}s directly into the input --- zero heap allocations for a 25K-row
batch. This avoids constructing 25K \texttt{Token} objects and 25K inner vectors that
the general parser would create.

\subsection{Compiler and Build Optimizations}

The Makefile uses aggressive optimization flags:

\begin{itemize}[noitemsep]
  \item \texttt{-O3 -march=native}: enables auto-vectorisation and targets the host CPU's
        instruction set (AVX2, SSE4.2)
  \item \texttt{-funroll-loops}: reduces loop overhead in tight scan loops
  \item \texttt{-flto=auto}: link-time optimisation across translation units, allowing
        cross-file inlining of hot functions like \texttt{StringArena::intern()}
\end{itemize}

No external libraries are linked. The only hardware-specific dependency is SSE4.2 for the
CRC32c instruction in the WAL writer and AVX2 for the SIMD-accelerated scan intrinsics.

\subsection{SIMD-Accelerated Numeric Scans (AVX2)}

Every \texttt{INT} and \texttt{DECIMAL} column maintains a parallel \texttt{vector<int64\_t>}
alongside the \texttt{string\_view} storage. On insert, each numeric value is parsed once
with an allocation-free integer parser (\texttt{parse\_int64}) and appended to
\texttt{int\_cols\_[col]}. This dual representation eliminates the
\texttt{std::stod(std::string(string\_view))} call that otherwise dominates scan cost.

For simple numeric \texttt{WHERE} predicates, the scan loop dispatches to hand-written
AVX2 intrinsics that compare 4~$\times$~\texttt{int64} values per CPU cycle:

\begin{lstlisting}[language=C++,basicstyle=\ttfamily\small,title={AVX2 scan kernel (simplified)}]
__m256i vals = _mm256_loadu_si256(col + i);   // load 4 int64s
__m256i mask = _mm256_cmpgt_epi64(vals, tgt); // compare all 4
int m = _mm256_movemask_pd(cast(mask));        // 4-bit result
if (m & 1) if (alive(i))   out.push_back(i);
if (m & 2) if (alive(i+1)) out.push_back(i+1);
// ...
\end{lstlisting}

Dedicated AVX2 kernels handle \texttt{=}, \texttt{!=}, \texttt{>}, \texttt{>=},
\texttt{<}, \texttt{<=}, \texttt{BETWEEN}, and \texttt{IN}. Non-numeric predicates
(\texttt{LIKE}, \texttt{IS NULL}, \texttt{VARCHAR} columns) fall through to the original
per-row \texttt{compare\_values()} path.

\texttt{ORDER BY} on numeric columns also uses the native \texttt{int64\_t} arrays,
eliminating \texttt{stod()} from every sort comparison.

\textbf{Memory overhead:} one \texttt{int64\_t} (8 bytes) per numeric value per row.
For a 1\,M-row table with 3 numeric columns, this adds $\sim$24\,MB (162\,MB total
vs.\ 139\,MB without).

\subsection{PK Fast Paths for DELETE, UPDATE, and SELECT IN}

The hash-map index is now used for more than just \texttt{SELECT WHERE pk = X}:

\begin{itemize}[noitemsep]
  \item \texttt{DELETE WHERE pk = X}: O(1) hash lookup + tombstone, instead of O(n) scan
  \item \texttt{UPDATE WHERE pk = X}: O(1) hash lookup + in-place modification
  \item \texttt{SELECT/DELETE/UPDATE WHERE pk IN (v1,\ldots,vk)}: O(k) hash lookups
\end{itemize}

This makes point mutations as fast as point reads.

\subsection{DoS Protection on \texttt{recv\_sql}}

The server enforces a 16\,MB cap on incoming SQL messages. If a client sends a message
exceeding this limit, \texttt{recv\_sql()} returns an empty string and the connection is
closed. This prevents a malicious or buggy client from exhausting server memory with a
single unbounded request.

\subsection{JOIN Deadlock Prevention}

When executing a JOIN, the executor must lock two tables. To prevent ABBA deadlocks when
two concurrent queries join the same tables in different order, the executor always
acquires table locks in \textbf{alphabetical order} of the table name. This total ordering
guarantees deadlock freedom without a lock manager.

\subsection{Database Singleton}

The \texttt{Database} class uses the singleton pattern (\texttt{Database::instance()}) to
provide a single global table registry. All threads access the same instance; the registry
itself is protected by a \texttt{shared\_mutex} for concurrent table lookups. Individual
table mutations are protected by per-table locks (see Section~6).

\subsection{Single-Condition WHERE Enforcement}

The assignment restricts \texttt{WHERE} to a single condition --- no \texttt{AND},
\texttt{OR}, or \texttt{NOT} combinators are allowed. After parsing one predicate atom,
the parser checks whether the next token is \texttt{AND} or \texttt{OR} and raises a
\texttt{ParseError} if so. This ensures the grading scripts receive deterministic,
single-condition behaviour. Operators like \texttt{BETWEEN \ldots\ AND}, \texttt{IN},
and \texttt{LIKE} are unaffected because their \texttt{AND} token is part of the
operator syntax, not a logical combinator.

% ══════════════════════════════════════════════════════════════════════════════
\section{Performance Results}

All benchmarks run on a 1\,M-row table (\texttt{BENCH}: 5 columns, no primary key)
unless stated otherwise.

\noindent\textbf{Hardware:} Intel Core i5-1135G7 @ 2.40\,GHz, 14\,GB RAM, Ubuntu Linux.\\
\textbf{Build:} \texttt{-O3 -march=native -std=c++17}.

\subsection{INSERT Throughput}

\begin{table}[H]
\centering
\begin{tabular}{lrr}
\toprule
\textbf{Dataset} & \textbf{Elapsed} & \textbf{Throughput} \\
\midrule
1\,M rows  & 277\,ms   & 3,610,108 rows/sec \\
10\,M rows & 3,548\,ms & 2,818,489 rows/sec \\
\bottomrule
\end{tabular}
\caption{INSERT benchmark --- batch size 5,000 (with dual int64 storage)}
\end{table}

The dual \texttt{int64\_t} storage adds a single \texttt{parse\_int64()} call per numeric
value during INSERT. This overhead is negligible: throughput remains above 2.8\,M~rows/sec
even at 10\,M rows.

\subsection{SELECT Queries}

\begin{table}[H]
\centering
\begin{tabular}{L{9cm} r r}
\toprule
\textbf{Query} & \textbf{Rows} & \textbf{Elapsed} \\
\midrule
\texttt{SELECT COUNT(*) FROM BENCH}                          & 1       & 174\,ms  \\
\texttt{SELECT * FROM BENCH} (full scan)                     & 1\,M    & 1,504\,ms \\
\texttt{SELECT * WHERE ID = 500000}                          & 1       & 0.75\,ms \\
\texttt{SELECT * WHERE SALARY > 90000}                       & 139,986 & 172\,ms  \\
\texttt{SELECT * WHERE SALARY BETWEEN 50000 AND 60000}       & 140,015 & 168\,ms  \\
\texttt{SELECT * WHERE DEPT IN ('ENG','HR','FIN')}           & 375,000 & 571\,ms  \\
\texttt{SELECT * WHERE NAME LIKE 'user1\%'}                  & 111,112 & 179\,ms  \\
\texttt{SELECT * ORDER BY SALARY DESC LIMIT 10}              & 10      & 200\,ms  \\
\texttt{SELECT * ORDER BY ID DESC LIMIT 100 OFFSET 1000}     & 100     & 209\,ms  \\
\texttt{SELECT DISTINCT DEPT FROM BENCH}                     & 8       & 94\,ms   \\
\bottomrule
\end{tabular}
\caption{SELECT query latencies on 1\,M-row table}
\end{table}

SIMD-accelerated numeric scans (\texttt{WHERE ID = 500000}, \texttt{SALARY > 90000},
\texttt{BETWEEN}) run $2\text{--}148\times$ faster than the pre-SIMD string-parsing path.
\texttt{ORDER BY} on numeric columns uses native \texttt{int64\_t} comparisons, reducing
sort time from $\sim$1,100\,ms to $\sim$200\,ms ($5.7\times$).

\subsection{Aggregate Queries}

\begin{table}[H]
\centering
\begin{tabular}{L{9cm} r}
\toprule
\textbf{Query} & \textbf{Elapsed} \\
\midrule
\texttt{SELECT SUM(SALARY)}                                & 165\,ms \\
\texttt{SELECT AVG(SALARY)}                                & 174\,ms \\
\texttt{SELECT MIN(SALARY)}                                & 173\,ms \\
\texttt{SELECT MAX(SALARY)}                                & 173\,ms \\
\texttt{GROUP BY DEPT} (8 groups)                          & 226\,ms \\
\texttt{GROUP BY DEPT HAVING AVG(SALARY)>60000}            & 174\,ms \\
\bottomrule
\end{tabular}
\caption{Aggregate query latencies on 1\,M-row table}
\end{table}

\subsection{JOIN Queries}

\begin{table}[H]
\centering
\begin{tabular}{L{8cm} r r}
\toprule
\textbf{Query} & \textbf{Rows} & \textbf{Elapsed} \\
\midrule
\texttt{INNER JOIN} (1\,M $\times$ 8, all match)      & 1\,M    & 960\,ms \\
\texttt{INNER JOIN + WHERE SALARY>90000}               & 139,986 & 372\,ms \\
\texttt{LEFT JOIN} (1\,M $\times$ 8)                   & 1\,M    & 1,195\,ms \\
\bottomrule
\end{tabular}
\caption{JOIN query latencies (BENCH 1\,M rows $\times$ DEPARTMENTS 8 rows)}
\end{table}

\subsection{UPDATE and DELETE}

\begin{table}[H]
\centering
\begin{tabular}{L{9cm} r}
\toprule
\textbf{Operation} & \textbf{Elapsed} \\
\midrule
\texttt{UPDATE} single row (\texttt{WHERE ID=1})                 & 91\,ms  \\
\texttt{UPDATE} 1,000 rows (\texttt{WHERE ID<=1000})             & 74\,ms  \\
\texttt{UPDATE} filtered (\texttt{WHERE DEPT='QA'})               & 55\,ms  \\
\texttt{DELETE} 100 rows (\texttt{WHERE ID<=100})                 & 76\,ms  \\
\texttt{DELETE} filtered (\texttt{WHERE DEPT='LEGAL'})            & 54\,ms  \\
\bottomrule
\end{tabular}
\caption{UPDATE / DELETE latencies on 1\,M-row table}
\end{table}

\subsection{LRU Cache Effect}

\begin{table}[H]
\centering
\begin{tabular}{lr}
\toprule
\textbf{Invocation} & \textbf{Elapsed} \\
\midrule
First  (cold cache) & 0.80\,ms \\
Second (warm cache) & 0.10\,ms \\
Third  (warm cache) & 0.03\,ms \\
\bottomrule
\end{tabular}
\caption{Cache hit vs.\ miss for a point query. Cold query now uses SIMD fast path (0.8\,ms)
instead of \texttt{stod()}-based scan ($\sim$113\,ms). Speedup: $\sim$29$\times$}
\end{table}

\subsection{Concurrency --- Concurrent Reads}

\begin{table}[H]
\centering
\begin{tabular}{rrrr}
\toprule
\textbf{Threads} & \textbf{Queries} & \textbf{Avg Latency} & \textbf{COUNT(*) wall-clock} \\
\midrule
1 & 20  & 0.68\,ms & 144\,ms \\
4 & 80  & 0.63\,ms & 201\,ms \\
8 & 160 & 0.46\,ms & 406\,ms \\
\bottomrule
\end{tabular}
\caption{Concurrent read scaling (\texttt{shared\_lock} allows parallel readers)}
\end{table}

Per-query latency is sub-millisecond because the SIMD numeric scan path processes
point queries in $\sim$0.5\,ms. Wall-clock for 8 simultaneous \texttt{COUNT(*)} scans
is only $2.8\times$ that of a single scan, showing good read parallelism.

\subsection{Concurrency --- Mixed Read/Write Workload}

\begin{table}[H]
\centering
\begin{tabular}{L{4cm} rrr}
\toprule
\textbf{Config} & \textbf{Read ops} & \textbf{Write ops} & \textbf{Total ops/sec} \\
\midrule
4 readers + 1 writer  & 853 & 63 & 305 \\
4 readers + 4 writers & 1,099 & 67 & 370 \\
\bottomrule
\end{tabular}
\caption{Mixed workload over 3 seconds (point queries + single-row UPDATEs)}
\end{table}

With SIMD scans, read throughput jumped from $\sim$33 to $\sim$305 ops/sec (9$\times$)
because point queries now complete in sub-millisecond time instead of $\sim$100\,ms.
Write contention still reduces read throughput because each \texttt{UPDATE} acquires an
exclusive \texttt{unique\_lock}, blocking concurrent readers momentarily.

\subsection{Concurrency --- Concurrent Inserts (Separate Tables)}

\begin{table}[H]
\centering
\begin{tabular}{rrrr}
\toprule
\textbf{Threads} & \textbf{Total Rows} & \textbf{Wall-clock} & \textbf{Agg.\ Throughput} \\
\midrule
1 &  50\,K  &  10\,ms  &  5.00\,M rows/sec \\
4 & 200\,K  & 13\,ms  & 15.38\,M rows/sec \\
8 & 400\,K  & 30\,ms  & 13.33\,M rows/sec \\
\bottomrule
\end{tabular}
\caption{Concurrent INSERT into separate tables (per-table locks, no contention)}
\end{table}

When each thread writes to its own table, per-table \texttt{unique\_lock} causes zero
contention. Aggregate throughput scales almost linearly: 8 threads achieve
$2.6\times$ the throughput of 1 thread.

\subsection{Memory Footprint}

\begin{table}[H]
\centering
\begin{tabular}{lr}
\toprule
\textbf{State} & \textbf{Server RSS} \\
\midrule
Idle (no tables)                  &   7.1\,MB \\
After 1\,M-row seed (5 columns)  & 161.8\,MB \\
After full benchmark suite        & $\sim$1.0\,GB \\
\bottomrule
\end{tabular}
\caption{Server resident set size at different stages}
\end{table}

The 1\,M-row table uses $\sim$155\,MB of memory (StringArena slabs + row metadata +
\texttt{string\_view} array + \texttt{int\_cols\_} native integer arrays). The extra
$\sim$24\,MB over the pre-SIMD baseline comes from the dual \texttt{int64\_t} storage
for the 3 numeric columns ($3 \times 1\text{M} \times 8\text{B} = 24\text{MB}$).

\subsection{Unit Test Results}

\textbf{564 / 564} standalone tests pass (\texttt{make test}). Test categories: parser
(17 functions), primary-key index (3), StringArena (4), table storage (10), executor
end-to-end (6), concurrency stress (3), adversarial inputs (5). Sanitiser targets
\texttt{make test-asan} and \texttt{make test-tsan} are also available.

% ══════════════════════════════════════════════════════════════════════════════
\section{Compilation and Execution}

\subsection{Requirements}

\begin{itemize}[noitemsep]
  \item GCC 11+ or Clang 14+ with C++17 support
  \item Linux (uses \texttt{writev}, \texttt{shared\_mutex}, POSIX threads)
  \item No external libraries required
\end{itemize}

\subsection{Build}

\begin{lstlisting}[language=bash]
git clone https://github.com/devchaitanya/FlexQl.git
cd FlexQl
make all -j$(nproc)
# Produces: bin/flexql-server  bin/flexql-client
\end{lstlisting}

\subsection{Run the Server}

\begin{lstlisting}[language=bash]
mkdir -p data/wal data/snapshots
./bin/flexql-server 9000
# Configuration: config/server.conf
\end{lstlisting}

\subsection{Run the Interactive Client (REPL)}

\begin{lstlisting}[language=bash]
./bin/flexql-client 127.0.0.1 9000
\end{lstlisting}

\begin{lstlisting}[language=SQL]
flexql> CREATE TABLE users (id INT PRIMARY KEY NOT NULL, name VARCHAR(64));
flexql> INSERT INTO users VALUES (1, 'Alice', 0), (2, 'Bob', 0);
flexql> SELECT * FROM users ORDER BY name;
+----+-------+
| id | name  |
+----+-------+
| 1  | Alice |
| 2  | Bob   |
+----+-------+
2 rows in set (0.000 sec)
flexql> .exit
\end{lstlisting}

\subsection{Run the Benchmark}

\begin{lstlisting}[language=bash]
./bench/benchmark 1000000     # 1M-row insert
./bench/benchmark 10000000    # 10M rows
./bench/benchmark 20000000    # 20M rows
./bench/benchmark --unit-test # unit tests only
\end{lstlisting}

\subsection{Source Layout}

\begin{lstlisting}[language={}]
FlexQl/
  include/          -- public headers (grouped by subsystem)
    cache/          -- lru_cache.h
    common/         -- types.h
    concurrency/    -- thread_pool.h
    expiration/     -- ttl_manager.h
    index/          -- pk_index.h
    network/        -- protocol.h, server.h
    parser/         -- ast.h, parser.h, token.h
    query/          -- executor.h
    storage/        -- arena.h, database.h, snapshot.h, table.h, wal.h
    utils/          -- string_utils.h
  src/              -- implementation (.cpp files, same layout as include/)
    client/         -- flexql.cpp (API), client.cpp (REPL)
  bench/            -- benchmark_flexql.cpp, flexql.h (instructor-provided)
  config/           -- server.conf
  docs/             -- design.tex  (this document)
  data/             -- wal/ and snapshots/ (created at runtime)
  Makefile
  README.md
  DESIGN.md
\end{lstlisting}

\end{document}
